% from ega3-1-fr.tex
\oldpage[0\textsubscript{III}]{12}
$s\mapsto\eta_X(Y)$ transformant tout élément en l'élément neutre. On
vérifie alors que le fait que $\eta_X(Y)$ soit élément neutre de
$\gamma_X(Y)$ pour tout $Y\in C$ équivaut à dire que le morphisme composé
\[
  X\xrightarrow{(1_X,1_X)}X\times X
  \xrightarrow{1_X\times e_X}X\times X\xrightarrow{c_X}X,
\]
et l'analogue où on permute $1_X$ et $e_X$, sont tous deux \emph{égaux} à
$1_X$.
\end{env}

\begin{env}[8.2.8]
\label{0.8.2.8-fr}
On pourrait bien entendu multiplier sans peine les exemples de structures
algébriques dans les catégories. L'exemple des groupes a été traité avec
assez de détails, mais par la suite nous laisserons généralement au lecteur
le soin de développer des considérations analogues dans les exemples de
structures algébriques que nous rencontrerons.
\end{env}

\section{Ensembles constructibles}
\label{section:0.9-fr}

\subsection{Ensembles constructibles}
\label{subsection:0.9.1-fr}

\begin{definition}[9.1.1]
\label{0.9.1.1-fr}
On dit qu'une application continue $f:X\to Y$ est \emph{quasi-compacte} si
pour tout ouvert quasi-compact $U$ de $Y$, $f^{-1}(U)$ est quasi-compact. On
dit qu'une partie $Z$ d'un espace topologique $X$ est \emph{rétrocompacte
dans $X$} si l'injection canonique $Z\to X$ est quasi-compacte, autrement dit
si pour tout ouvert quasi-compact $U$ de $X$, $U\cap Z$ est quasi-compact.
\end{definition}

Une partie \emph{fermée} de $X$ est rétrocompacte dans $X$, mais une partie
quasi-compacte de $X$ n'est pas nécessairement rétrocompacte dans $X$. Si
$X$ est quasi-compact, toute partie ouverte rétrocompacte dans $X$ est
quasi-compacte. Il est clair que toute réunion \emph{finie} d'ensembles
rétrocompacts dans $X$ est rétrocompacte dans $X$, toute réunion finie
d'ensembles quasi-compacts étant quasi-compacte. Toute intersection finie
d'ouverts rétrocompacts dans $X$ est un ouvert rétrocompact dans $X$. Dans un
espace \emph{localement noethérien} $X$, tout ensemble quasi-compact est un
sous-espace noethérien, et par suite \emph{toute partie} de $X$ est
rétrocompacte dans $X$.

\begin{definition}[9.1.2]
\label{0.9.1.2-fr}
Étant donné un espace topologique $X$, on dit qu'une partie de $X$ est
\emph{constructible} si elle appartient au plus petit ensemble de parties
$\mathfrak{F}$ de $X$ contenant toutes les parties ouvertes rétrocompactes de
$\mathfrak{F}$ et stable par intersection finie et passage au complémentaire
(ce qui implique que $\mathfrak{F}$ est aussi stable par réunion finie).
\end{definition}

\begin{proposition}[9.1.3]
\label{0.9.1.3-fr}
Pour qu'une partie de $X$ soit constructible, il faut et il suffit qu'elle
soit réunion finie d'ensembles de la forme $U\cap\complement{V}$, où $U$ et
$V$ sont des ouverts rétrocompacts dans $X$.
\end{proposition}

Il est clair que la condition est suffisante. Pour voir qu'elle est
nécessaire, considérons l'ensemble $\mathfrak{G}$ des réunions finies
d'ensembles de la forme $U\cap\complement{V}$ où $U$ et $V$ sont ouverts
rétrocompacts dans $X$~; il suffit de voir que tout complémentaire d'un
ensemble de $\mathfrak{G}$ appartient à $\mathfrak{G}$. Soit donc
$Z=\bigcup_{i\in I}(U_i\cap\complement{V_i})$, où $I$ est fini, $U_i$ et
$V_i$ ouverts rétrocompacts dans $X$~; on a
$\complement{Z}=\bigcap_{i\in I}(V_i\cup\complement{U_i})$, donc
$\complement{Z}$ est réunion finie d'ensembles qui sont intersections d'un
certain nombre des $V_i$ et d'un certain nombre de $\complement{U_i}$, donc
de la forme
